\documentclass[11pt]{article} \usepackage[margin=1in]{geometry} \usepackage{hyperref} \usepackage{enumitem} \usepackage{tabularx} \usepackage{longtable} \usepackage{array} \usepackage{amsmath} \usepackage{adjustbox} \usepackage{pdflscape} \usepackage{makecell} \title{Warehouse Operations Planning: Final Report} \author{Amonios Beshara, Hazem Nasr, Mikhael Khalil, and Yehia El Hodeiby} \date{May 16, 2026} \begin{document} \maketitle \section*{Title and Team Roles} \textbf{Roles:} \begin{itemize}[leftmargin=1.25em] \item Thread A (Human): Yehia (enhanced algorithm design), Hazem (implementation and complexity analysis) \item Thread B (AI-assisted): Amonios (prompting and reasoning-audit lead), Mikhael (AI-assisted coding, testing, and verification) \item Integration and editing: Hazem (repo formatting and organization), Amonios (report integration and editing) \end{itemize} \textbf{Contribution Matrix:} \begin{table}[h] \centering \begin{tabularx}{\textwidth}{|l|c|c|c|c|c|} \hline Task & Weight (\%) & Amonios (\%) & Mikhael (\%) & Hazem (\%) & Yehia (\%) \\ \hline Problem modeling & 10 & 15 & 15 & 30 & 40 \\ Algorithm design (Thread A) & 25 & 0 & 0 & 35 & 65 \\ AI-assisted workflow (Thread B) & 20 & 60 & 40 & 0 & 0 \\ Implementation (C++ code) & 20 & 10 & 35 & 40 & 15 \\ Testing and outputs & 15 & 20 & 40 & 25 & 15 \\ Report integration and editing & 10 & 40 & 10 & 40 & 10 \\ \hline \end{tabularx} \caption{Contribution matrix.} \end{table} \textbf{Note on contributions.} Much of the work was done locally while sitting together, so the repo history does not fully reflect effort. Hazem had less algorithm-design work than Yehia but was responsible for formatting and organizing the repo to balance the workload. \section*{Formal Problem Definition} We must construct a feasible next-day warehouse operations plan. Inputs include a warehouse graph (dispatch node D, optional supplier node V, shelves), items with stock and handling constraints, customer orders with release times, deadlines, priorities, and packing times, picker shifts and capacity limits, and optional restock constraints (reserve, prep time, delivery time, and max delivery weight). The plan must assign orders to picker batches, schedule any needed restocks, and route each batch through shelves to meet deadlines and shift constraints while respecting stock feasibility, perishable handling limits, and routing feasibility. If the plan is feasible, output restock requests and batches with times and routes; otherwise report INFEASIBLE with a specific reason. \section*{Thread A: Human Algorithm Development} \textbf{Problem model.} The team modeled the warehouse as a weighted undirected graph and used all-pairs shortest paths for routing. Orders are batched per picker and routed by a nearest-neighbor heuristic. Feasibility is enforced through stock, capacity, release, shift, deadline, and perishable-time checks. \textbf{Algorithm overview.} The human thread implemented two greedy strategies: \begin{itemize}[leftmargin=1.25em] \item \textbf{Baseline greedy:} orders are sorted by earliest deadline and restocking is triggered only when stock runs out. \item \textbf{Enhanced greedy:} orders are scored by a weighted linear function of normalized features (priority, slack, stock urgency, expected handling time, and price). Restocks are scheduled with lead time, including buffers for high delay risk. \end{itemize} Batching uses a composite rule that selects the next order by maximizing $\text{score} - \lambda \cdot \text{normalizedDistance}$ with $\lambda = 0.2$, and routes are built by a nearest-neighbor tour over required shelves. \textbf{Correctness (feasibility) sketch.} Each batch is accepted only if stock is available (after scheduled restocks), weight capacity is satisfied, release and shift windows are respected, deadlines are met, and perishable max-wait is not exceeded. Therefore, any output plan is feasible by construction when the algorithm succeeds. \textbf{Complexity (high-level).} Let $V$ be the number of nodes, $M$ the number of orders, $P$ the number of pickers, and $S$ the number of shelves per batch. The dominant preprocessing cost is Floyd--Warshall $O(V^3)$. Sorting orders is $O(M\log M)$, batch construction is $O(M^2)$, and routing is $O(S^2)$ per batch. \textbf{Thread A test results.} \begin{table}[h] \centering \begin{tabular}{|c|c|c|c|c|c|} \hline Case & Baseline Feasible & Enhanced Feasible & Baseline Restocks & Enhanced Restocks & Runtime (ms) \\ \hline 1 & Yes & Yes & 0 & 0 & 18.58 \\ 2 & Yes & Yes & 3 & 2 & 13.58 \\ 3 & No (unassigned) & Yes & 0 & 8 & 39.13 \\ 4 & No (VIP deadline) & No (VIP deadline) & 0 & 0 & 4.60 \\ \hline \end{tabular} \caption{Thread A test results for cases 1--4.} \end{table} \section*{Thread B: AI-Assisted Development} \textbf{AI workflow summary.} The AI-assisted thread implemented a priority-driven greedy planner with restock scheduling, batching, and feasibility validation. The algorithm phases are: (1) all-pairs shortest paths, (2) demand analysis, (3) restock scheduling with delay buffers, (4) order prioritization by priority and slack, (5) capacity-aware greedy batching, (6) greedy routing via nearest-neighbor TSP, and (7) feasibility validation. \textbf{Complexity (from code comments).} Floyd--Warshall $O(V^3)$, sorting $O(M\log M)$, batching $O(M\cdot P)$, routing $O(S^2)$ per batch, and validation $O(M\cdot S)$. Overall $O(V^3 + M\cdot P + B\cdot S^2)$ where $B$ is the number of batches. \textbf{Thread B execution outputs.} \begin{table}[h] \centering \begin{tabular}{|c|c|c|c|} \hline Case & Feasible & Restock Requests & Runtime (ms) \\ \hline 1 & Yes & 0 & 6.32 \\ 2 & Yes & 2 & 10.17 \\ 3 & Yes & 8 & 34.19 \\ 4 & No (VIP deadline) & [N/A] & 4.98 \\ \hline \end{tabular} \caption{Thread B execution summary for cases 1--4.} \end{table} \textbf{Thread B test notes.} Runs were deterministic (fixed seed = 42) and all four cases executed on the same machine within the 60-second limit. During development, cases 1--3 initially returned INFEASIBLE due to a batch-packing bug, which was fixed by a round-robin batch loop and restock-arrival guards. Manual edge cases tested included: zero-demand order (empty lines), single picker with one fully stocked order, exact reserve equals shortfall, reserve less than shortfall (INFEASIBLE), and overtime trigger by shrinking shift end time. \textbf{Human edits and corrections.} We revised AI output in four main places: (1) restock splitting by max delivery weight, (2) priority-aware order sorting (VIP and urgent tiers with deadline tie-breaks), (3) perishable shelf-wait checks, and (4) round-robin batch packing so orders were not assigned to a single picker. These edits fixed shallow scheduling assumptions and reduced false confidence in feasibility. \begin{landscape} \section*{Good Prompts, Weak Prompts, and Bad Prompts} \textbf{Discussion lead:} Thread B. \begin{longtable}{|c|p{3cm}|c|p{4cm}|p{3cm}|c|c|p{2cm}|} \hline Prompt \# & Prompt summary & Good/Weak/Bad & Why? & AI output & Understanding? & Productivity? & Risk \\ \hline 1 & Initial framing: design algorithm from vague prompt & Bad & Too broad; missing constraints & Generic greedy scheduler & No & Medium & High \\ \hline 2 & Graph model and Floyd-Warshall complexity & Good & Specific graph and request & Correct F-W and $O(V^3)$ & Yes & Medium & Low \\ \hline 3 & Order sorting to minimize lateness & Weak & Missing priority tiers & EDF only, incomplete & Partial & Medium & Medium \\ \hline 4 & Restock scheduling code (no context) & Bad & Missing constraints & Stub with three bugs & No & Medium & High \\ \hline 5 & Multi-delivery split with max weight & Good & Exact constraints given & Correct splitting loop & Yes & Medium & Low \\ \hline 6 & Nearest-neighbor TSP heuristic & Good & Clear start/stop/return & Correct NN, $O(S^2)$ & Yes & Medium & Low \\ \hline 7 & Perishable handling question & Weak & Too vague & Generic expiry advice & No & Low & Medium \\ \hline 8 & Batch packing bug with code snippet & Good & Specific bug and code & Round-robin fix & Yes & Medium & Low \\ \hline 9 & Infeasibility reporting & Good & Clear failure modes & Structured reasons & Yes & Medium & Low \\ \hline 10 & Case 2 stock violation debug & Good & Failure time and code & Guard on restock arrival & Yes & Medium & Low \\ \hline \caption{Required prompt quality table.} \end{longtable} \end{landscape} \textbf{Prompt quality summary.} We documented 12 interactions: 5 good, 3 weak, and 4 bad prompts. The good prompts always included constraints or code context; the bad prompts asked for generation without context and produced plausible but wrong output. \textbf{AI as reviewer vs. generator.} The most reliable use of AI was as a reviewer after we had a design. When we asked it to generate without context, it introduced subtle bugs that took more time to fix than writing the logic ourselves. \section*{Getting the Answer Right vs. Understanding the Answer} We evaluate understanding using the following checks (complete at least four): \begin{itemize}[leftmargin=1.25em] \item Explain the algorithm without reading AI output: Achieved for Thread A; partial for Thread B. \item Trace the algorithm on a new example: Achieved with one new small case per thread. \item Derive time complexity step-by-step: Achieved for both threads. \item Explain why the greedy choice is valid: Partial; relied on empirical justification. \item Identify an edge case where a naive solution fails: Achieved (VIP deadline case). \item Modify the algorithm under a new constraint: Achieved (restock max weight constraint). \item Explain why the C++ implementation matches the pseudocode: Achieved after review. \item Answer oral questions about the algorithm and code: Achieved for main components. \end{itemize} Completion sentence: \begin{quote} AI helped us get an answer in a fast, structured way, but understanding required manual tracing, constraint checks, and corrections. \end{quote} \section*{Productivity Comparison} \textbf{Time logs.} \begin{table}[h] \centering \small \setlength{\tabcolsep}{4pt} \begin{tabularx}{\textwidth}{|l|l|X|l|l|X|} \hline Student & Date & Task & AI used? & Time spent & Output produced \\ \hline Yehia & 2026-04-22 & Enhanced algorithm design & No & 6 h & Scoring model and batch rule \\ \hline Hazem & 2026-04-23 & C++ implementation and complexity & Yes & 5 h & Thread A code and analysis \\ \hline Mikhael & 2026-04-24 & AI-assisted coding and tests & Yes & 4 h & Thread B code and run logs \\ \hline Amonios & 2026-04-25 & Prompt audit and report writing & Yes & 4 h & Prompt log and analysis notes \\ \hline \end{tabularx} \caption{Productivity time log.} \end{table} \textbf{Summary metrics.} \begin{table}[h] \centering \begin{tabularx}{\textwidth}{|l|c|c|X|} \hline Measurement & Human thread & AI-assisted thread & Interpretation \\ \hline Time to first problem model & 60 min & 30 min & AI sped drafting but missed constraints \\ \hline Time to first algorithm idea & 90 min & 40 min & Faster with AI, weaker justification \\ \hline Time to first pseudocode & 2.5 h & 1.5 h & AI accelerated structure \\ \hline Time to first working C++ code & 6 h & 3 h & AI sped coding, required review \\ \hline Number of major revisions & 2 & 4 & AI output needed more correction \\ \hline Number of edge cases discovered & 3 & 4 & AI suggested more cases \\ \hline Number of incorrect claims found & 1 & 3 & AI reasoning was shallow \\ \hline Number of failed tests & 1 & 3 & Early AI batches failed feasibility \\ \hline Time spent debugging or correcting & 2 h & 3.5 h & AI saved time early, cost time later \\ \hline Final confidence in correctness & High & Medium & Confidence tied to manual verification \\ \hline \end{tabularx} \caption{Productivity comparison summary.} \end{table} \section*{Correctness and Complexity Comparison} \textbf{Correctness evidence.} Thread A and Thread B both produced feasible plans for cases 1--3 and identified the infeasible VIP deadline in case 4. Thread A reported a baseline failure in case 3 while the enhanced variant succeeded. Thread B required manual fixes for restock splitting, priority ordering, perishable checks, and round-robin batching to avoid incorrect feasibility. \textbf{Complexity comparison.} \begin{table}[h] \centering \begin{tabular}{|l|l|} \hline Thread & Stated Complexity \\ \hline Thread A (Human) & $O(V^3 + M^2 + S^2)$ with $O(V^3)$ dominating \\ Thread B (AI-assisted) & $O(V^3 + M\cdot P + B\cdot S^2)$ \\ \hline \end{tabular} \caption{High-level complexity comparison.} \end{table} \textbf{Interpretation.} Both approaches share a $O(V^3)$ preprocessing bottleneck. Thread A's enhanced scoring can push toward $O(M^2)$, while Thread B trades for $O(M\cdot P)$ but depends more on validity checks. Complexity claims were similar, but correctness depended on careful constraint handling. \section*{When AI Helped (Concrete Examples)} Provide at least three examples from the group work. For each, state the impact. \begin{enumerate} \item AI suggested a phased planner structure that matched the problem constraints.\newline Impact: Clarity and productivity. \item AI generated edge-case ideas (VIP deadline infeasible case) that we adopted.\newline Impact: Testing and correctness. \item AI accelerated C++ scaffolding for data structures and output formatting.\newline Impact: Productivity. \end{enumerate} \section*{When AI Harmed or Could Have Harmed (Concrete Examples)} Provide at least three examples from the group work. \begin{enumerate} \item AI initially ignored restock weight limits, which would allow infeasible deliveries. \item AI provided a shallow correctness argument that masked missing constraints. \item AI suggested batching by capacity only, which risks deadline violations. \item AI missed perishable window checks until prompted with exact constraints. \end{enumerate} \section*{Recommendation for Future CSCE 2202 Students} Use AI after an initial human attempt. Start with manual problem modeling and at least one baseline algorithm, then use AI to review constraints, generate edge cases, and scaffold code. Require manual verification of correctness and complexity claims before accepting any AI-generated reasoning. \section*{Final Message on AI Usage and the Core Project Question} AI helps students and professionals work faster when it is used as a reviewer or testing assistant after the core reasoning is done. It weakens algorithmic skills when it substitutes for modeling, proof thinking, or manual tracing. In academic settings, the cost of shallow understanding is higher because those skills are the learning outcomes; in professional settings, speed gains are acceptable only if verification practices are strong. \appendix \section*{Appendix A: Source Code and Execution Logs} \begin{itemize}[leftmargin=1.25em] \item Human solution code: \url{https://github.com/HazeMOONium/Analysis-Project/blob/main/human_solution.cpp} (local file: \texttt{\detokenize{human_solution.cpp}}) \item AI-assisted solution code: \url{https://github.com/HazeMOONium/Analysis-Project/blob/main/ai_solution.cpp} (local file: \texttt{\detokenize{ai_solution.cpp}}) \item Human execution output: \url{https://github.com/HazeMOONium/Analysis-Project/tree/main/human_output} (local file: \texttt{\detokenize{human_output/output.txt}}) \item AI execution output cases 1--4: \url{https://github.com/HazeMOONium/Analysis-Project/tree/main/ai_output} (local files: \texttt{\detokenize{ai_output/execution_output_case_*.txt}}) \end{itemize} \section*{Appendix B: AI Usage Appendix} \begin{itemize}[leftmargin=1.25em] \item AI tool used: GitHub Copilot Chat (GPT-5.2-Codex) \item Dates of AI use: 2026-05-14 to 2026-05-16 \item Prompts: Summarized in the 10-row prompt log above \item AI outputs: Algorithm outline, C++ scaffolding, test ideas, debugging suggestions \item Accepted by group: Phase structure, NN TSP, infeasibility reporting \item Rejected by group: EDF-only sorting, restock stubs without max weight, vague perishable advice \item Corrected by group: Restock splitting, priority tiers, perishable constraint, round-robin batching \end{itemize} \end{document}